% Clase 3 — Planeación: Heurísticas y Planeación Jerárquica
% Lectura: AIMA Secciones 11.3–11.4
\section{Planeación: Heurísticas y Planeación Jerárquica}

\subsection{Heurísticas para Planeación}

\subsubsection{Problema Relajado (ignore preconditions / ignore delete lists)}
\begin{itemize}
    \item \textbf{Ignore preconditions}: todas las acciones son aplicables. Heurística: número mínimo de acciones para satisfacer la meta.
    \item \textbf{Ignore delete lists}: no se eliminan fluentes. Problema de planeación relajado resoluble en tiempo polinomial.
\end{itemize}

\subsubsection{Planning Graph y Heurística FF}
Un \textbf{grafo de planeación} alterna capas de literales y acciones:
\begin{itemize}
    \item Capa de literales $S_i$: literales alcanzables en $i$ pasos.
    \item Capa de acciones $A_i$: acciones aplicables en el paso $i$.
\end{itemize}
La heurística FF usa el plan del grafo relajado como estimador.

\subsection{Planeación Jerárquica (HTN)}

\textbf{Hierarchical Task Network (HTN)}: descompone tareas abstractas en subtareas más concretas.

\begin{itemize}
    \item \textbf{Tareas primitivas}: acciones directamente ejecutables.
    \item \textbf{Tareas abstractas}: se descomponen mediante métodos.
    \item Ventaja: permite incorporar conocimiento de dominio para guiar la búsqueda.
\end{itemize}
